\section{Anexos}\label{sec:anexos}

\subsection{Estructura de carpetas del proyecto}\label{sec:anexo-estructura}

\begin{lstlisting}[basicstyle=\small\ttfamily]
TFM-ML-Benchmark/
|-- data/
|   |-- raw/
|   `-- processed/support2/
|       |-- df1/  # estrategia clinico-estadistica
|       `-- df2/  # estrategia tecnico-estadistica
|-- eda/
|-- preprocessing/
|-- models/
|-- scripts/modeling_utils_tidymodels_clean.R
|-- figures/
`-- manuscript/
\end{lstlisting}

\subsection{Diccionario de variables del dataset SUPPORT2}\label{sec:anexo-diccionario}

\begingroup
\small
\setlength{\tabcolsep}{4pt}
\begin{longtable}{p{2.1cm}p{1.4cm}p{6.7cm}p{3.8cm}}
\caption{Diccionario de variables del dataset SUPPORT2.}\label{tab:diccionario-support2}\\
\toprule
\textbf{Variable} & \textbf{Tipo} & \textbf{Descripción} & \textbf{Unidades / Valores posibles} \\
\midrule
\endfirsthead
\toprule
\textbf{Variable} & \textbf{Tipo} & \textbf{Descripción} & \textbf{Unidades / Valores posibles} \\
\midrule
\endhead
id & Entero & Identificador único del paciente en el estudio. & Número único \\
age & Continuo & Edad del paciente al ingresar al estudio. & Años \\
death & Continuo & Fallecimiento en cualquier momento del seguimiento (hasta el 31/DIC/1994). & 0 = No, 1 = Sí \\
sex & Categórico & Género biológico del paciente. & male (hombre), female (mujer) \\
hospdead & Binario & Fallecimiento ocurrido dentro del hospital durante el ingreso. & 0 = No, 1 = Sí \\
slos & Continuo & Días transcurridos desde el ingreso al estudio hasta el alta. & Días \\
d.time & Continuo & Días totales de seguimiento del paciente en el estudio. & Días \\
dzgroup & Categórico & Subcategoría específica del diagnóstico principal de ingreso. & ARF/MOSF con Sepsis, CHF, COPD, Cirrosis, Cáncer de Colon, Coma, Cáncer de Pulmón, MOSF con Malignidad \\
dzclass & Categórico & Categoría general de la enfermedad (agrupación de dzgroup). & ``ARF/MOSF'', ``COPD/CHF/Cirrhosis'', ``Cancer'', ``Coma'' \\
num.co & Continuo & Número de comorbilidades (enfermedades simultáneas). A mayor valor, peor pronóstico. & Número ordinal \\
edu & Categórico & Nivel educativo medido en años de escolaridad. & Años \\
income & Categórico & Nivel estimado de ingresos económicos anuales del paciente. & ``under \$11k'', ``\$11-\$25k'', ``\$25-\$50k'', ``>\$50k'' \\
scoma & Continuo & Puntuación de coma basada en la escala de Glasgow en el día 3 (predicha por modelo). & Escala numérica \\
charges & Continuo & Cargos o costos económicos facturados por el hospital. & Dólares (\$) \\
totcst & Continuo & Costo total calculado mediante la relación costo-cargo (RCC). & Dólares (\$) \\
totmcst & Continuo & Micro-costo total calculado para el paciente. & Dólares (\$) \\
avtisst & Continuo & Puntuación promedio TISS (días 3-25). Mide la intensidad de intervenciones terapéuticas en UCI. & Escala TISS \\
race & Categórico & Etnia o raza registrada del paciente. & asian, black, hispanic, missing, other, white \\
sps & Continuo & Puntuación de fisiología SUPPORT en el día 3 (predicha por modelo). & Escala numérica \\
aps & Continuo & Puntuación fisiológica APACHE III en el día 3 (sin coma, con BUN corregido y diuresis). & Escala numérica \\
surv2m & Continuo & Estimación del modelo SUPPORT de la probabilidad de supervivencia a los 2 meses (en el día 3). & Probabilidad (0 a 1) \\
surv6m & Continuo & Estimación del modelo SUPPORT de la probabilidad de supervivencia a los 6 meses (en el día 3). & Probabilidad (0 a 1) \\
hday & Entero & Día de hospitalización en el que el paciente ingresó formalmente al estudio. & Días \\
diabetes & Continuo & Presencia de diabetes mellitus como enfermedad comórbida preexistente. & 0 = No, 1 = Sí \\
dementia & Continuo & Presencia de demencia crónica como enfermedad comórbida preexistente. & 0 = No, 1 = Sí \\
ca & Categórico & Estado de enfermedad oncológica activa en el paciente. & no (sano), yes (cáncer localizado), metastatic (metástasis) \\
prg2m & Continuo & Estimación subjetiva del médico sobre la probabilidad de supervivencia a 2 meses. & Probabilidad (0 a 1) \\
prg6m & Categórico & Estimación subjetiva del médico sobre la probabilidad de supervivencia a 6 meses. & Probabilidad (0 a 1) \\
dnr & Categórico & Estado de la orden de No Reanimación (DNR) y el momento de su firma. & no dnr, dnr before sadm, dnr after sadm, missing \\
dnrday & Continuo & Día del ingreso en que se firmó la orden DNR (valores negativos indican pre-ingreso). & Días \\
meanbp & Continuo & Presión arterial media (PAM) del paciente registrada en el día 3. & mmHg \\
wblc & Continuo & Conteo total de glóbulos blancos en sangre en el día 3. & Miles por $\mu$L \\
hrt & Continuo & Frecuencia cardíaca registrada en el día 3. & Latidos por minuto (lpm) \\
resp & Continuo & Frecuencia respiratoria registrada en el día 3. & Respiraciones por minuto (rpm) \\
temp & Continuo & Temperatura corporal medida en el día 3. & Grados Celsius (°C) \\
pafi & Continuo & Índice de oxigenación Kirby ($PaO_2/FiO_2$). Indicador clínico clave de hipoxemia. & Ratio (mmHg) \\
alb & Continuo & Nivel de albúmina sérica (proteína) medido en el día 3. & g/dL \\
bili & Continuo & Nivel de bilirrubina total en sangre medido en el día 3. & mg/dL \\
crea & Continuo & Nivel de creatinina sérica en sangre (función renal) medido en el día 3. & mg/dL \\
sod & Continuo & Concentración de sodio sérico en sangre medida en el día 3. & mEq/L \\
ph & Continuo & pH de la sangre arterial medido en el día 3 (equilibrio ácido-base). & Escala de pH (6.8 a 7.8) \\
glucose & Entero & Nivel de glucosa en sangre medido en el día 3. & mg/dL \\
bun & Entero & Nitrógeno ureico en sangre (función renal/metabólica) medido en el día 3. & mg/dL \\
urine & Entero & Gasto urinario (diuresis total de 24 horas) medido en el día 3. & mL \\
adlp & Categórico & Índice de Actividades de la Vida Diaria (ADL) respondido por el propio paciente. & Escala ordinal \\
adls & Continuo & Índice de Actividades de la Vida Diaria (ADL) respondido por un familiar/representante. & Escala ordinal \\
sfdm2 & Categórico & Escala de discapacidad funcional (SIP) a los 2 meses. El nivel 5 denota muerte temprana. & Escala 1 al 5 \\
adlsc & Continuo & Índice ADL imputado y calibrado estadísticamente con respecto al reporte del familiar. & Escala numérica \\
\bottomrule
\end{longtable}
\endgroup
